\section{Adversarial Coherence Failures}
\label{sec:adversarial}

The naturalistic experiments in §\ref{sec:paradox} demonstrate the
coherence paradox \emph{in distribution}, but leave open the question of
whether our coherence signals can isolate \emph{specific} fragmentation
patterns. This section reports an adversarial probe with hand-authored
cases that target three failure modes individually.

\subsection{Case design}

We construct 40 cases organized into four subsets of 10 (data files in
\texttt{data/adversarial/}; loader at \texttt{src/adversarial\_cases.py}).
Each case is a tuple \texttt{(case\_id, category, question, passages,
expected\_failure\_mode, notes)}. The categories are:

\begin{description}
  \item[\textsc{Disjoint}] passages that are individually relevant to the
        question but drawn from disjoint topical neighborhoods. Example:
        ``\emph{What is the primary cause of myocardial infarction?}''
        with passages alternating between cardiology, ICU triage, and
        sports-medicine recovery — every passage matches a query keyword
        but none cohere.
  \item[\textsc{Contradict}] relevant passages containing contradictory
        claims about the queried entity (e.g., one passage says aspirin
        reduces stroke risk, another says it raises bleeding risk in the
        same cohort). Faithfulness here cannot be measured against a
        single ground-truth answer; the targeted signal is the
        contradiction itself.
  \item[\textsc{Drift}] passages where the topic gradually drifts from a
        near-miss neighbor to the true target (``capital of Australia''
        passages drifting Sydney → Canberra → ACT museums). The model
        must locate the answer in only one of the passages.
  \item[\textsc{Control}] matched coherent controls: passages that share
        the same surface topic and length distribution but do not exhibit
        any of the three pathologies. These calibrate signal magnitudes
        against a non-adversarial baseline.
\end{description}

The cases are released as part of ContextCoherenceBench
(§\ref{sec:benchmark}) so other researchers can probe their own
retrievers.

\subsection{Detection task}

For each case we compute eight retrieval-time signals:

\begin{enumerate}
  \item \texttt{mean\_query\_chunk\_sim} (negated)
  \item \texttt{ccs} (negated)
  \item \texttt{mean\_jaccard} (negated)
  \item \texttt{embedding\_variance}
  \item \texttt{retrieval\_entropy}
  \item \texttt{sim\_spread}
  \item \texttt{nli\_pairwise\_mean}
  \item \texttt{nli\_pairwise\_max}
  \item \texttt{nli\_pairwise\_frac\_hi}
\end{enumerate}

Signals (1)--(3) are negated so that, for every signal, ``higher = more
adversarial'' under the predicted direction. We then report binary
detection AUC (sklearn \texttt{roc\_auc\_score}) for each (signal,
adversarial-category-vs-control) pair.

The \texttt{nli\_pairwise\_*} family is new in this paper:
\texttt{compute\_nli\_pairwise} (in \texttt{src/coherence\_metrics.py})
runs the same NLI cross-encoder used for faithfulness scoring on every
pair of retrieved passages and reports the mean / max contradiction
score and the fraction of pairs above a contradiction threshold. This
gives the \textsc{Contradict} category a signal that pure embedding-space
metrics cannot capture by construction.

\subsection{Preregistered analysis plan}
\label{sec:adv_plan}

As with the mechanistic probe, we publish the analysis plan before the
numbers so that the scoring cannot drift to fit the signal panel. The
final reporting unit is a single table of detection AUCs: three rows
(one per adversarial category: \textsc{disjoint}, \textsc{contradict},
\textsc{drift}) and eight columns (one per coherence signal listed
above). Each cell is the AUC for separating the row's adversarial subset
from the matched \textsc{control} subset using that column's signal;
cells at AUC $\geq 0.80$ are marked in bold. We also report 95\%
bootstrap confidence intervals over the 10 paired items per row so a
reader can tell high-but-noisy cells apart from high-and-stable ones.

We commit in advance to two success conditions: (i) at least one
coherence signal achieves AUC $\geq 0.80$ on each of the three
adversarial categories, and (ii) the best signal for \textsc{contradict}
comes from the \texttt{nli\_pairwise\_*} family rather than from any
embedding-space signal, which is the dissociation the category was
designed to produce. A uniform AUC table — whether uniformly high,
uniformly low, or uniformly middling — counts as a failed dissociation
and will be reported as written.

\subsection{What we expect to learn}

The categories were designed to dissociate signals:

\begin{itemize}
  \item \textsc{Disjoint} should be detected best by
        \texttt{ccs}, \texttt{mean\_jaccard}, and \texttt{embedding\_variance}
        — embedding-space coherence is exactly what fragments under
        topical disjointness.
  \item \textsc{Contradict} should be detected by
        \texttt{nli\_pairwise\_*} and only weakly by embedding signals,
        because contradictory passages can be \emph{embedded} close
        together (similar topic, opposite claim).
  \item \textsc{Drift} should be detected by \texttt{sim\_spread} and
        \texttt{retrieval\_entropy} — drift produces a long tail of
        decreasing similarity rather than a homogeneous low-similarity
        set.
\end{itemize}

If this dissociation holds empirically, the practical implication is that
no single coherence signal is sufficient: a deployment monitor wanting to
catch all three failure modes needs at least one embedding-space signal
\emph{and} an NLI-pair signal. We therefore recommend treating the
signals as a panel rather than a scalar.

\paragraph{Honest negative case.} If the AUC table shows uniformly high
or uniformly low values across categories and signals, the probe has
failed to distinguish failure modes — either the cases are too easy / too
hard, or our signal panel is collinear in practice. We will report this
outcome as written if it occurs; the probe is informative under both
results.
